% !TeX root = ../main.tex
\section{Benchmarking results}
\label{sec:benchmarking-results}

\newcommand*{\micros}[1]{\num[round-precision=2,round-mode=places]{\fpeval{1000000*(#1)}}}
\newcommand*{\millis}[1]{\num[round-precision=2,round-mode=places]{\fpeval{1000*(#1)}}}

\begin{figure*}
    \includegraphics[width=0.7\linewidth]{bh-total-time.pdf-legend.pdf} \\[-5pt]
    \begin{subfigure}[b]{0.32\textwidth}
        \includegraphics[width=\linewidth]{bh-total-time.pdf}
        \vspace{-20pt}
        \subcaption{} \label{fig:bh-total-time}
    \end{subfigure}
    \begin{subfigure}[b]{0.32\textwidth}
        \includegraphics[width=\linewidth]{bh-sample-time.pdf}
        \vspace{-20pt}
        \subcaption{} \label{fig:bh-sample-time}
    \end{subfigure}
    \begin{subfigure}[b]{0.32\textwidth}
        \includegraphics[width=\linewidth]{bh-memory.pdf}
        \vspace{-20pt}
        \subcaption{} \label{fig:bh-memory}
    \end{subfigure}
    \caption{
        Comparison of different simulators for 1000 attempted shots of the Bravyi-Haah $3k+8$ to $k$ magic state distillation protocol \cite{bravyi2012magic} for increasing numbers of logical qubits $k$.
        (\subref{fig:bh-total-time})~Total runtime, including pre-processing.
        (\subref{fig:bh-sample-time})~Average time per shot, excluding pre-processing.
        (\subref{fig:bh-memory})~Peak memory consumption.
    }
    \label{fig:bh}
\end{figure*}

\setlength{\tabcolsep}{5pt}
\begin{table}[t]
    \centering
    \textbf{Total runtime (including pre-processing) in \qty{}{\milli\second}}
    \begin{tabular}{lrrrrr}
        \toprule
        \textbf{Circuit} &
        \texttt{tsim} &
        \texttt{xtim} &
        \texttt{clifft} &
        \texttt{symft} &
        \texttt{merlin} \\
        \midrule
        15-to-1  & \millis{5.029947541002912} & \millis{0.08585524999944028} & \millis{0.09955558299770928} & \millis{0.057665624997753184} & \textbf{\millis{0.012521416996605694}} \\
        Cultiv-3 & OOM & \millis{0.06529049999517156} & \millis{0.05429354199804948} & \millis{0.06362920799801941} & \textbf{\millis{0.02898095799901057}} \\
        Cultiv-5 & OOM & \millis{0.2535097500003758} & \textbf{\millis{0.067934000002424}} & \millis{0.5772002910016454} & \millis{0.20201220999661018} \\
        Switch-3 & OOM & \millis{2.817529874999309} & OOM & ERR & \textbf{\millis{0.5736940430033428}} \\
        Switch-5 & OOM & OOM & OOM & ERR & \textbf{\millis{6.959127001005982}} \\
        \bottomrule
    \end{tabular}

    \bigskip
    \textbf{Time per shot (excluding pre-processing) in \qty{}{\micro\second}}
    \begin{tabular}{lrrrrr}
        \toprule
        \textbf{Circuit} &
        \textbf{tsim} &
        \textbf{xtim} &
        \textbf{clifft} &
        \textbf{symft} &
        \textbf{merlin} \\
        \midrule
        15-to-1  & \micros{0.0034468530410013045} & \micros{4.025839989481028e-07} & \micros{3.707920004671905e-07} & \textbf{\micros{1.8404200091026724e-07}} & \micros{9.620625001844019e-06} \\
        Cultiv-3 & OOM & \micros{3.898329996445682e-07} & \micros{5.06666001456324e-07} & \textbf{\micros{3.4187499841209503e-07}} & \micros{2.8341749999526654e-05} \\
        Cultiv-5 & OOM & \micros{7.704665997152915e-06} & \micros{7.97404200056917e-06} & \textbf{\micros{6.517500001791632e-06}} & \micros{0.00020014883399926476} \\
        Switch-3 & OOM & \textbf{\micros{2.652375002071494e-06}} & OOM & ERR & \micros{0.0005499323750009353} \\
        Switch-5 & OOM & OOM & OOM & ERR & \textbf{\micros{0.006704935917001421}} \\
        \bottomrule
    \end{tabular}

    \medskip
    \raggedright
    \footnotesize
    \quad OOM: Simulator tried to allocate more than 20GB of memory. \\
    \quad ERR: Simulator failed with a runtime error.

    \caption{
        Simulation time for 1000 attempted shots of the 15-to-1 magic state distillation \cite{bravyi2005universal}, magic state cultivation without the escape stage \cite{gidney2024magic} and the code switching protocol from Ref. \cite{li2026transversal}.
    }
    \label{tab:other-circuits}
\end{table}

We provide an open-source implementation of our simulator, which we call \texttt{merlin} \cite{koch2026merlin}.
In this section, we benchmark \texttt{merlin} against existing non-Clifford simulation libraries on a range of magic-state preparation circuits.
Concretely, we compare against Tsim~\cite{haenel2026tsim}, Clifft~\cite{chase2026clifft}, SymFT~\cite{fang2026symft}, and Xtim~\cite{surti2026efficient}.
All benchmarks use a noise parameter $p=10^{-3}$ and 1000 attempted shots per circuit.
The cultivation benchmarks exclude the escape stage.

While these simulators each use different representations and algorithms, they are all particularly effective when the number of non-stabilizer degrees of freedom remains small.
This is a favorable regime for many existing magic-state preparation protocols.
At the logical level, the target states often only represent a small number of magic states $k$, with $k = 1$ being especially simple.
These simulators exploit the fact that this simplicity can persist after the protocol is expanded into a noisy physical circuit.
However, their cost can increase rapidly as the amount of logical magic grows.
In contrast, our approach does not require $k$ to remain fixed or small.
Our state representation and update rules have polynomial complexity in the physical system size and circuit size, and therefore introduce no exponential dependence on the number of logical magic states $k$.

Figure~\ref{fig:bh} illustrates this difference using the Bravyi-Haah $3k+8$ to $k$ magic state distillation protocol \cite{bravyi2012magic} for increasing numbers of logical output magic states $k$.
The other simulators exhibit exponentially increasing runtime and/or memory consumption as $k$ grows, consistent with the growth of the non-stabilizer representations that they maintain.
In contrast, the runtime and memory consumption of \texttt{merlin} grow smoothly with $k$, reflecting the polynomial scaling of our representation.

Table~\ref{tab:other-circuits} reports results for the 15-to-1 magic state distillation~\cite{bravyi2005universal}, magic-state cultivation~\cite{gidney2024magic} and the code switching protocol for $CCZ$ magic states introduced in Ref. \cite{li2026transversal}.
The distillation and cultivation circuits have previously been used as benchmark cases for non-Clifford simulators~\cite{chase2026clifft,fang2026symft,surti2026efficient}.
They each prepare a single logical magic state and therefore lie in the $k=1$ regime, where Xtim, Clifft, and SymFT achieve substantially lower per-shot runtimes than \texttt{merlin}.
For the code switching benchmarks, we use the $[\![27,3,3]\!]$ and $[\![81,3,5]\!]$ BT code constructions of \cite{li2026transversal} which both encode $k = 3$ logical qubits.
On these circuits, the other simulators' techniques to keep their non-stabilizer representations compact appear to be less effective.
Only Xtim completes the distance-3 simulation, while all other simulator-circuit combinations run out of memory or fail with a runtime error.

These examples highlight the complementary regime targeted by our approach:
while specialized methods can have smaller constant factors when the effective non-stabilizer representation remains small, \texttt{merlin} can handle circuits for which such representations cannot be found or become prohibitively large.
